\documentclass[12pt]{article}
\usepackage[utf8]{inputenc}
\usepackage[T1]{fontenc}
\usepackage[margin=1in]{geometry}
\usepackage{setspace}
\setlength{\parindent}{2em}
\setlength{\parskip}{1.25em}
\renewcommand{\baselinestretch}{1.0}
\makeatletter
\newcommand{\@subtitle}{}
\newcommand{\subtitle}[1]{\renewcommand{\@subtitle}{#1}}
\renewcommand{\@maketitle}{%
  \newpage
  \null
  \vskip 1em%
  \begin{center}%
    {\LARGE \@title \par}%
    \ifx\@subtitle\@empty\else
      \vskip 0.5em%
      {\large \@subtitle \par}%
    \fi
  \end{center}%
  \par
  \vskip 1em}
\makeatother
\usepackage{xcolor}
\usepackage[colorlinks=true,urlcolor=blue!55!black,linkcolor=black]{hyperref}

\title{\textit{Dred Scott v. Sandford}, 60 U.S. 393 (1857)}
\subtitle{Case Brief}
\author{}
\date{}

\begin{document}
\maketitle

\section{Memory Jogger}

In \href{https://www.courtlistener.com/opinion/87116/dred-scott-v-sandford/}{\textit{Dred Scott v. Sandford}}, the Supreme Court held that people of African descent could not be citizens of the United States under the Constitution and struck down the Missouri Compromise of 1820, accelerating the sectional crisis that led to the Civil War before being repudiated by the Thirteenth and Fourteenth Amendments.

\section{Detailed Case Facts}

Dred Scott was born into slavery in Virginia around 1799 and was later brought by his enslaver, Peter Blow, to St.\ Louis, Missouri, where Blow sold him to U.S. Army surgeon Dr.\ John Emerson (\href{https://en.wikipedia.org/wiki/Dred_Scott_v._Sandford}{Wikipedia}; \href{https://www.sos.mo.gov/archives/resources/africanamerican/scott/scott.asp}{Missouri Digital Heritage}). In 1834, Dr.\ Emerson took Scott from Missouri (a slave state) to the military post at Rock Island in Illinois, a free state where slavery had been prohibited by the Northwest Ordinance of 1787 and the Illinois Constitution of 1819, and held him there until April or May 1836 (\href{https://msa.maryland.gov/megafile/msa/speccol/sc5300/sc5339/000091/000000/000004/restricted/dred_scott/decision/nelson.html}{Maryland State Archives}). Emerson then took Scott to Fort Snelling, located on the west bank of the Mississippi River in the Wisconsin Territory (part of the former Louisiana Purchase north of latitude $36^{\circ}30'$ north, in present-day Minnesota), where Congress had ``forever prohibited'' slavery under Section 8 of the Missouri Compromise of 1820 (Act of March 6, 1820, 3 Stat.\ 545) (\href{https://www.archives.gov/milestone-documents/dred-scott-v-sandford}{National Archives}).

At Fort Snelling in 1836, Scott married Harriet Robinson, an enslaved woman brought there by Army officer and Indian agent Major Lawrence Taliaferro and transferred to Dr.\ Emerson (\href{https://msa.maryland.gov/megafile/msa/speccol/sc5300/sc5339/000091/000000/000004/restricted/dred_scott/decision/daniel.html}{Maryland State Archives}). When reassigned to Jefferson Barracks near St.\ Louis and later Fort Jesup in Louisiana, Emerson initially left Dred and Harriet Scott behind at Fort Snelling and hired out their labor for profit (\href{https://en.wikipedia.org/wiki/Dred_Scott_v._Sandford}{Wikipedia}). After Emerson married Eliza Irene Sanford in Louisiana in 1838, he summoned the Scotts and subsequently returned with them to Fort Snelling; during the steamboat journey north of Missouri on the Mississippi River, their daughter Eliza was born, and a second daughter, Lizzie, was later born in Missouri (\href{https://msa.maryland.gov/megafile/msa/speccol/sc5300/sc5339/000091/000000/000004/restricted/dred_scott/decision/daniel.html}{Maryland State Archives}). After returning to St.\ Louis, Dr.\ Emerson died in 1843, leaving his estate to his widow Irene, who continued to hire out the Scotts and refused Dred Scott's attempt in 1846 to purchase his family's freedom (\href{https://en.wikipedia.org/wiki/Dred_Scott_v._Sandford}{Wikipedia}).

The litigation implicated three major constitutional provisions: Article III, Section 2 (diversity jurisdiction over controversies ``between Citizens of different States''); Article IV, Section 3, Clause 2 (the Territory or Property Clause, empowering Congress ``to dispose of and make all needful Rules and Regulations respecting the Territory or other Property belonging to the United States''); and the Fifth Amendment's Due Process Clause (\href{https://supreme.justia.com/cases/federal/us/60/393/}{Justia}). Politically, the case unfolded amid explosive national conflict over the expansion of slavery into western territories following the Kansas--Nebraska Act of 1854 and ``Bleeding Kansas.'' Behind the scenes, President-elect James Buchanan improperly corresponded with Justices John Catron and Robert C. Grier, pressuring Grier (a Pennsylvanian) to join the Southern majority so the ruling would appear national rather than sectional, and declaring in his March 4, 1857 inaugural address that the slavery question would be ``speedily and finally settled'' by the Court (\href{https://en.wikipedia.org/wiki/Dred_Scott_v._Sandford}{Wikipedia}).

\section{Procedural History}

\textbf{Missouri State Courts (\textit{Scott v. Emerson}):} On April 6, 1846, Dred and Harriet Scott filed freedom suits in the Circuit Court of St.\ Louis County against Irene Emerson, relying on settled Missouri precedent (including \textit{Winny v. Whitesides}, 1 Mo.\ 472 (1824), and \textit{Rachel v. Walker}, 4 Mo.\ 350 (1836)) establishing ``once free, always free'' for enslaved persons taken to reside in free jurisdictions (\href{https://en.wikipedia.org/wiki/Dred_Scott_v._Sandford}{Wikipedia}). After losing an initial June 1847 trial on a hearsay ruling, Scott won a new trial; in January 1850, a jury returned a verdict freeing Scott and Judge Alexander Hamilton declared Dred, Harriet, Eliza, and Lizzie free (\href{https://en.wikipedia.org/wiki/Dred_Scott_v._Sandford}{Wikipedia}). On March 22, 1852, however, a politically reconstituted Missouri Supreme Court reversed 2--1 in \textit{Scott v. Emerson}, 15 Mo.\ 576 (1852), repudiating its prior precedents on states'-rights and comity grounds and holding that Scott remained a slave upon returning to Missouri (\href{https://www.sos.mo.gov/archives/resources/africanamerican/scott/scott.asp}{Missouri Digital Heritage}).

\textbf{U.S. Circuit Court for the District of Missouri (\textit{Scott v. Sanford}):} With the state case stalled on remand and Irene Emerson having moved to Massachusetts, Scott's counsel Roswell Field filed a new diversity action of trespass \textit{vi et armis} in federal court in 1853 against Irene's brother, John F. A. Sanford, a New York citizen who had assumed control over the Scotts (the Supreme Court clerk later misspelled his surname as ``Sandford'') (\href{https://www.oyez.org/cases/1850-1900/60us393}{Oyez}; \href{https://en.wikipedia.org/wiki/Dred_Scott_v._Sandford}{Wikipedia}). Sanford filed a plea in abatement challenging federal diversity jurisdiction on the ground that Scott was a ``negro of African descent'' whose ancestors were imported and sold as slaves, and therefore not a ``citizen'' of Missouri (\href{https://supreme.justia.com/cases/federal/us/60/393/}{Justia}). Judge Robert W. Wells sustained Scott's demurrer to the plea in abatement in April 1854, ruling that Scott could sue if free, and Sanford pleaded over to the merits (\href{https://supreme.justia.com/cases/federal/us/60/393/}{Justia}). On an agreed statement of facts in May 1854, Judge Wells instructed the jury that Missouri law governed and favored Sanford; the jury returned a verdict for Sanford (\href{https://supreme.justia.com/cases/federal/us/60/393/}{Justia}).

\textbf{U.S. Supreme Court:} Scott brought the case to the Supreme Court by writ of error. Represented by Montgomery Blair and George Ticknor Curtis (brother of Justice Benjamin R. Curtis), against Reverdy Johnson and Henry S. Geyer for Sanford, the case was argued February 11--14, 1856, and reargued December 15--18, 1856 (\href{https://supreme.justia.com/cases/federal/us/60/393/}{Justia}; \href{https://www.oyez.org/cases/1850-1900/60us393}{Oyez}). On March 6, 1857, the Court issued a 7--2 decision reversing the Circuit Court's judgment for Sanford and remanding with directions to dismiss the suit for want of jurisdiction (\href{https://www.archives.gov/milestone-documents/dred-scott-v-sandford}{National Archives}).

\section{Judicial Votes}

The Court ruled \textbf{7--2} against Dred Scott, with \textbf{all nine Justices filing separate (seriatim-style) opinions} totaling over 200 pages in the \textit{United States Reports} (\href{https://tile.loc.gov/storage-services/service/ll/usrep/usrep060/usrep060393a/usrep060393a.pdf}{Library of Congress}; \href{https://constitutioncenter.org/the-constitution/supreme-court-case-library/dred-scott-v-sandford}{National Constitution Center}). Although Chief Justice Taney's opinion was designated the ``Opinion of the Court,'' the seven-justice majority fractured over which issues were properly before the Court:

\begin{itemize}
  \item \textbf{Opinion of the Court (1 author, joined in full by 1):} Chief Justice Roger B. Taney (author), joined without qualification by Justice James M. Wayne (who also wrote a short concurrence).
  \item \textbf{Separate Concurrences (6):}
    \begin{itemize}
      \item Justice James M. Wayne (concurred completely in Taney's opinion on all points).
      \item Justice Samuel Nelson (concurred in the result on Missouri law alone under \textit{Strader v. Graham}; declined to reach the plea in abatement or the Missouri Compromise; would have affirmed).
      \item Justice Robert C. Grier (joined Nelson's opinion on Missouri law, and also concurred with Taney that the Missouri Compromise was void and that Scott could not sue as a citizen).
      \item Justice Peter V. Daniel (concurred with Taney on the plea in abatement, Black non-citizenship, and the unconstitutionality of the Missouri Compromise).
      \item Justice John A. Campbell (concurred in the judgment; declined to rule on the plea in abatement, relied on Missouri law, and held the Missouri Compromise unconstitutional under the Territory Clause).
      \item Justice John Catron (concurred in the result; held the plea in abatement was waived, joined Nelson on Illinois/Missouri law, affirmed general congressional power over territories, but invalidated the Missouri Compromise under the 1803 Louisiana Purchase Treaty and state equality; would have affirmed).
    \end{itemize}
  \item \textbf{Separate Dissents (2):} Justice John McLean (dissented in full); Justice Benjamin R. Curtis (dissented in full).
\end{itemize}

\section{Holding}

\textbf{Judgment:} Judgment of the U.S. Circuit Court reversed and case remanded with directions to dismiss the suit for lack of federal jurisdiction (\href{https://www.archives.gov/milestone-documents/dred-scott-v-sandford}{National Archives}).

\textbf{Rules Announced:} (1) Persons of African descent whose ancestors were imported into the United States and sold as slaves---whether enslaved or free---are not and cannot become ``citizens'' within the meaning of the U.S. Constitution and therefore cannot invoke federal diversity jurisdiction under Article III; (2) Section 8 of the Missouri Compromise of 1820, prohibiting slavery in federal territories north of $36^{\circ}30'$, is unconstitutional and void (6--3 across Taney, Wayne, Grier, Daniel, Campbell, and Catron) because Congress lacks power to exclude slave property from the territories and doing so deprives slaveholders of property without due process of law under the Fifth Amendment; and (3) under \textit{Strader v. Graham}, 51 U.S. (10 How.) 82 (1851), Scott's temporary residence in Illinois did not emancipate him once he returned to Missouri, whose courts had final authority over his status (\href{https://supreme.justia.com/cases/federal/us/60/393/}{Justia}; \href{https://www.archives.gov/milestone-documents/dred-scott-v-sandford}{National Archives}).

\textbf{Holding vs.\ Dicta:} The boundary between holding and \textit{dicta} in \textit{Dred Scott} became a central legal and political battleground. Critics---beginning with Justice Curtis in dissent and Abraham Lincoln in the 1858 debates---contended that once the Court ruled on the plea in abatement that it lacked Article III jurisdiction, every subsequent pronouncement on the Missouri Compromise and Fifth Amendment was extrajudicial \textit{obiter dictum} (\href{https://en.wikipedia.org/wiki/Dred_Scott_v._Sandford}{Wikipedia}). Taney and Wayne countered that because federal courts of limited jurisdiction must inspect the entire record on a writ of error from a Circuit Court, the merits of Scott's freedom claim independently bore on whether he was a citizen capable of invoking jurisdiction (\href{https://msa.maryland.gov/megafile/msa/speccol/sc5300/sc5339/000091/000000/000004/restricted/dred_scott/decision/wayne.html}{Maryland State Archives}). Moreover, because only three justices (Taney, Wayne, and Daniel) explicitly endorsed the plea-in-abatement holding excluding all free Black people from citizenship, while six justices (Taney, Wayne, Grier, Daniel, Campbell, and Catron) explicitly held the Missouri Compromise unconstitutional, the territorial ruling actually commanded a clearer numerical majority on the Court than the broad citizenship ruling.

\section{Analysis of Opinions}

\subsection{Opinion of the Court (Chief Justice Taney)}

Chief Justice Taney organized his opinion around two sweeping propositions. First, reviewing the plea in abatement, he held that no person of African descent whose ancestors were brought to America as slaves could ever be a ``citizen'' under the Constitution (\href{https://www.archives.gov/milestone-documents/dred-scott-v-sandford}{National Archives}). Taney distinguished state citizenship from federal citizenship, asserting that while a State could confer local rights, only the Constitution and Congress's uniform naturalization power defined U.S. citizenship. Appealing to his reading of 1776 and 1787 public opinion, colonial and state discriminatory statutes, the Slave Trade Clause, and the Fugitive Slave Clause, Taney asserted that Black people had been viewed as ``beings of an inferior order'' with ``no rights which the white man was bound to respect'' and were excluded from ``We the People'' and the Declaration of Independence (\href{https://constitutioncenter.org/the-constitution/supreme-court-case-library/dred-scott-v-sandford}{National Constitution Center}). He warned that recognizing Black citizenship would entitle free Black Americans to ``the full liberty of speech,'' public assembly, and the right ``to keep and carry arms'' nationwide under the Privileges and Immunities Clause---consequences he deemed unthinkable to the Framers (\href{https://en.wikipedia.org/wiki/Dred_Scott_v._Sandford}{Wikipedia}).

Second, reaching the merits to determine whether Scott's residence in free territory emancipated him, Taney invalidated the Missouri Compromise---the first federal statute struck down since \textit{Marbury v. Madison} (1803) (\href{https://en.wikipedia.org/wiki/Dred_Scott_v._Sandford}{Wikipedia}). He narrowly construed the Article IV Territory Clause as applying solely to territory already belonging to the United States in 1787 (such as the Northwest Territory), not to later acquisitions like the Louisiana Purchase (\href{https://supreme.justia.com/cases/federal/us/60/393/}{Justia}). Even where Congress governed new territories as a trustee for the States, it remained bound by the Bill of Rights. Invoking an early form of substantive due process under the Fifth Amendment, Taney declared that ``the right of property in a slave is distinctly and expressly affirmed in the Constitution'' and that an Act of Congress depriving a citizen of slave property merely for bringing it into a federal territory ``could hardly be dignified with the name of due process of law'' (\href{https://constitutioncenter.org/the-constitution/supreme-court-case-library/dred-scott-v-sandford}{National Constitution Center}). Finally, under \textit{Strader v. Graham}, Scott's stay in Illinois did not override Missouri law upon his return.

\subsection{Concurring Opinion (Justice Wayne)}

Justice Wayne filed a brief concurrence expressing ``unqualified assent'' to every line of Taney's opinion (\href{https://msa.maryland.gov/megafile/msa/speccol/sc5300/sc5339/000091/000000/000004/restricted/dred_scott/decision/wayne.html}{Maryland State Archives}). Wayne candidly explained that ``the peace and harmony of the country required the settlement'' of the slavery controversy by judicial decision, and defended the Court's authority to decide both jurisdiction and the Missouri Compromise on a writ of error from a federal Circuit Court.

\subsection{Concurring Opinion (Justice Nelson)}

Justice Nelson wrote the narrowest opinion in the majority (originally drafted to serve as the Court's opinion before Southern justices pushed for a broader ruling). Declining to rule on either the plea in abatement or the constitutionality of the Missouri Compromise, Nelson rested solely on conflict-of-laws principles and \textit{Strader v. Graham} (\href{https://msa.maryland.gov/megafile/msa/speccol/sc5300/sc5339/000091/000000/000004/restricted/dred_scott/decision/nelson.html}{Maryland State Archives}). Citing Lord Stowell's decision in \textit{The Slave Grace} (1827) and Justice Story's \textit{Conflict of Laws}, Nelson reasoned that because neither Illinois law nor congressional territorial legislation had extraterritorial force inside Missouri, Missouri's supreme authority over domestic status controlled once Scott returned, requiring affirmance of the Circuit Court's judgment.

\subsection{Concurring Opinion (Justice Grier)}

In a single-paragraph concurrence, Justice Grier joined Nelson's opinion on Missouri law while also concurring with Taney that the Missouri Compromise of 1820 was ``unconstitutional and void'' and that Scott could not sue as a citizen of Missouri (\href{https://msa.maryland.gov/megafile/msa/speccol/sc5300/sc5339/000091/000000/000004/restricted/dred_scott/decision/grier.html}{Maryland State Archives}). Grier remarked that because Scott was a slave on the merits, whether the judgment was affirmed or dismissed for want of jurisdiction was ``of little importance.''

\subsection{Concurring Opinion (Justice Daniel)}

Justice Daniel delivered an intensely proslavery concurrence agreeing with Taney on the plea in abatement, Black non-citizenship, and the invalidity of the Missouri Compromise (\href{https://msa.maryland.gov/megafile/msa/speccol/sc5300/sc5339/000091/000000/000004/restricted/dred_scott/decision/daniel.html}{Maryland State Archives}). Drawing on Roman law, Justinian's \textit{Institutes}, and Emer de Vattel's \textit{Law of Nations}, Daniel argued that enslaved persons were strictly private property incapable of being parties to the social compact, that individual manumission could never create citizenship, and that slave property was the only form of private property affirmatively guaranteed by the Constitution. He rejected the Northwest Ordinance of 1787 as void \textit{ab initio} and held the Territory Clause conferred only proprietary authority over public lands.

\subsection{Concurring Opinion (Justice Campbell)}

Justice Campbell declined to pass on the plea in abatement, focusing instead on the merits and the Territory Clause (\href{https://msa.maryland.gov/megafile/msa/speccol/sc5300/sc5339/000091/000000/000004/restricted/dred_scott/decision/campbell.html}{Maryland State Archives}). He concluded that Dr.\ Emerson's military postings effected no change of domicile, so Missouri law governed Scott's status upon return. Turning to the Missouri Compromise, Campbell argued from Revolutionary-era constitutional history and the Tenth Amendment that the Territory Clause authorized only administrative rules for surveying and selling the public domain, not municipal sovereignty to dissolve master--slave relations established by state law.

\subsection{Concurring Opinion (Justice Catron)}

Justice Catron broke sharply with Taney on two major points while still invalidating the Missouri Compromise (\href{https://msa.maryland.gov/megafile/msa/speccol/sc5300/sc5339/000091/000000/000004/restricted/dred_scott/decision/catron.html}{Maryland State Archives}). First, he held that the plea in abatement was waived when Sanford pleaded over to the merits and was not properly before the Court. Second, he rejected Taney's narrow reading of the Territory Clause, insisting that sixty years of practice and judicial precedent confirmed Congress's broad power to govern territories. Nevertheless, joining Nelson on Illinois law, Catron struck down the Missouri Compromise on two distinct grounds: (1) it violated Article III of the 1803 Louisiana Purchase Treaty, which protected inhabitants' property (including slaves), and (2) it violated the constitutional equality of the States under the Privileges and Immunities Clause by excluding slaveholding citizens from common national territory.

\subsection{Dissenting Opinion (Justice McLean)}

Justice McLean dissented across the board (\href{https://msa.maryland.gov/megafile/msa/speccol/sc5300/sc5339/000091/000000/000004/restricted/dred_scott/decision/McLean.html}{Maryland State Archives}; \href{https://constitutioncenter.org/the-constitution/supreme-court-case-library/dred-scott-v-sandford}{National Constitution Center}). On jurisdiction, he argued that citizenship for Article III purposes required only freedom and permanent domicile in a State, not the right to vote (noting that women and minors were citizens who could sue in federal court), and that free Black men had exercised suffrage in New England at the Founding. On the merits, quoting \textit{Prigg v. Pennsylvania}, McLean insisted that ``[t]he state of slavery is deemed to be a mere municipal regulation, founded upon and limited to the range of the territorial laws''; once Dr.\ Emerson voluntarily took Scott into free Illinois and Fort Snelling, slavery ceased to operate and could not be reimposed (\href{https://msa.maryland.gov/megafile/msa/speccol/sc5300/sc5339/000091/000000/000004/restricted/dred_scott/decision/McLean.html}{Maryland State Archives}). He defended Congress's plenary authority under the Territory Clause to prohibit slavery and attacked the majority's territorial ruling as unauthorized \textit{dicta} once it found no jurisdiction.

\subsection{Dissenting Opinion (Justice Curtis)}

In a meticulous 67-page dissent that prompted Taney to add 18 pages of rebuttal before publication, Justice Curtis dismantled both of Taney's core holdings (\href{https://en.wikipedia.org/wiki/Dred_Scott_v._Sandford}{Wikipedia}; \href{https://famous-trials.com/dredscott/2545-scott-v-sandford-the-dissent}{Famous Trials}). On citizenship, Curtis proved from historical records that at the ratification of the Constitution, free native-born Black men were not only citizens but qualified voters in five of the thirteen original States---New Hampshire, Massachusetts, New York, New Jersey, and North Carolina---and therefore belonged to ``the people of the United States'' who ordained the Constitution (\href{https://constitutioncenter.org/the-constitution/supreme-court-case-library/dred-scott-v-sandford}{National Constitution Center}). He further demonstrated that citizenship does not depend on possessing every political or civil right, pointing to women, minors, and D.C.\ residents. On the merits, Curtis argued that once the majority found no jurisdiction on the plea in abatement, its invalidation of the Missouri Compromise was wholly extrajudicial; on the substance, he defended Congress's broad power under the Territory Clause (evidenced by the Northwest Ordinance and decades of congressional legislation), refuted Taney's Fifth Amendment reading by noting that slavery was a creation of positive local law rather than natural property, and concluded that Scott's marriage and residence at Fort Snelling permanently emancipated him and his family.

\section{Examples: Future Applications}

\textit{Note:} Because \textit{Dred Scott} was constitutionally superseded by the Thirteenth Amendment (1865) and the Citizenship Clause of the Fourteenth Amendment (1868), the following hypotheticals illustrate how its doctrinal logic would have operated under the \textit{Dred Scott} regime itself, as well as how its methodological issues echo in modern constitutional disputes.

\begin{enumerate}
  \item \textbf{Same-Side Application (Free Black Merchant Suing in Federal Diversity Court in 1858):} A free-born Black merchant from Massachusetts sues a New York shipper in federal Circuit Court for breach of a maritime contract exceeding the amount-in-controversy threshold. Under Taney's holding in \textit{Dred Scott}, the federal court must dismiss the suit for lack of Article III diversity jurisdiction, because Massachusetts's recognition of the merchant as a state citizen cannot make him a ``citizen'' of the United States.
  \item \textbf{Same-Side Application (Territorial Legislature Banning Slavery in Kansas or Nebraska):} Following \textit{Dred Scott}, the elected territorial legislature of Nebraska enacts a statute banning slavery within the territory under ``popular sovereignty.'' Under Taney's Fifth Amendment reasoning---that Congress cannot delegate power it does not itself possess---the territorial ban would be unconstitutional and void because neither Congress nor its territorial creature may deprive a citizen of slave property brought into a federal territory.
  \item \textbf{Opposite-Side Application (State-Court Freedom Suit in a Northern Free State):} A Virginia slaveholder brings an enslaved person to reside in Massachusetts for two years, and the enslaved person files a writ of habeas corpus in Massachusetts state court while still physically in Massachusetts. Even under \textit{Dred Scott}, the Massachusetts court could grant freedom, because Taney and Nelson's holdings turned on \textit{Missouri}'s sovereignty after the enslaved person's return and did not (yet) hold that free States were constitutionally required to recognize slavery within their own borders.
  \item \textbf{Opposite-Side Application (Post-Fourteenth Amendment Birthright Citizenship Challenge):} A State attempts to deny state or federal citizenship to U.S.-born children of a disfavored racial minority by invoking historical exclusion at the Founding. Under the Fourteenth Amendment's Citizenship Clause---enacted specifically to overturn \textit{Dred Scott}---such a law is flatly unconstitutional, as confirmed in \textit{United States v. Wong Kim Ark}, 169 U.S. 649 (1898), which repudiated Taney's racial-exclusion framework.
  \item \textbf{Fence-Sitter Application (Slaveholder's Right of Transit Through a Free State---the \textit{Lemmon} Problem):} A slaveholder traveling from Virginia to Texas stops overnight in New York City with enslaved persons, and New York officials invoke a state statute freeing any enslaved person brought voluntarily into New York. This was the explosive ``next case'' Justice Nelson expressly reserved at the end of his \textit{Dred Scott} concurrence (\href{https://msa.maryland.gov/megafile/msa/speccol/sc5300/sc5339/000091/000000/000004/restricted/dred_scott/decision/nelson.html}{Maryland State Archives}): under \textit{Dred Scott}'s logic of state sovereignty over domestic status, New York's law should govern within New York; yet under Taney's and Catron's sweeping assertions of a constitutional right to carry slave property under the Privileges and Immunities and Due Process Clauses, Abraham Lincoln warned in his ``House Divided'' speech that a second \textit{Dred Scott} decision could have nationalized slavery even inside free States.
\end{enumerate}

\section{Critique}

\textbf{Universal Condemnation and Progressive / Historical Criticism:} Across virtually the entire modern legal spectrum, \textit{Dred Scott} occupies the apex of the constitutional ``anticanon,'' famously labeled by Chief Justice Charles Evans Hughes as the Court's ``greatest self-inflicted wound'' (\href{https://en.wikipedia.org/wiki/Dred_Scott_v._Sandford}{Wikipedia}). Historians such as Don E. Fehrenbacher and Paul Finkelman have documented that Taney's opinion combined overt white supremacist ideology with gross historical distortion (\href{https://studentorgs.kentlaw.iit.edu/cklawreview/wp-content/uploads/sites/3/vol82no1/Finkelman.pdf}{Finkelman, Chicago-Kent Law Review}). As Justice Curtis proved in real time, free Black men voted in five states in 1787--1788, directly refuting Taney's claim of universal Founding-era exclusion. Moreover, Taney's reading of the Territory Clause as limited to 1787 lands contradicted both its plain text and seventy years of uninterrupted congressional governance of the Louisiana and Florida territories. Politically, the Court's hubristic attempt to ``settle'' the slavery crisis by judicial fiat---amplified by President Buchanan's improper back-channel lobbying---destroyed the Court's credibility in the North, fractured the Democratic Party, and hastened the Civil War (\href{https://lawcat.berkeley.edu/record/1124652/files/fulltext.pdf}{Farber, Berkeley Law}).

\textbf{Originalist and Textualist Perspectives:} \textit{Dred Scott} is a focal point of intense debate over originalism and substantive due process (\href{https://www.law.georgetown.edu/public-policy-journal/wp-content/uploads/sites/23/2018/10/16-2-Dred-Scott-Roe.pdf}{Georgetown Journal of Law \& Public Policy}). Critics of originalism point out that Taney himself employed overtly originalist rhetoric---insisting that the Constitution ``must be construed and administered now according to its true meaning and intention when it was formed and adopted''---as a cautionary tale showing how easily original-intent arguments can be manipulated to entrench injustice (\href{https://www.archives.gov/milestone-documents/dred-scott-v-sandford}{National Archives}; \href{https://studentorgs.kentlaw.iit.edu/cklawreview/wp-content/uploads/sites/3/vol82no1/Finkelman.pdf}{Finkelman}). Modern originalists and textualists respond on two fronts. First, they point to Justice Curtis's dissent as the model of disciplined, text-and-history originalism, demonstrating that Taney got the Founding-era text and history demonstrably wrong (including claiming the Constitution ``expressly affirmed'' property in slaves when the text carefully referred to enslaved people as ``persons''). Second, conservative jurists---including Justice Antonin Scalia in \textit{Casey} and Justice Clarence Thomas in \textit{Dobbs}---have indicted \textit{Dred Scott} as the ``original precedent'' for substantive due process, arguing that Taney's invention of a substantive Fifth Amendment right to carry slave property into federal territories exemplifies the danger of judges using the Due Process Clause to override democratic legislation (\href{https://en.wikipedia.org/wiki/Dred_Scott_v._Sandford}{Wikipedia}; \href{https://www.law.georgetown.edu/public-policy-journal/wp-content/uploads/sites/23/2018/10/16-2-Dred-Scott-Roe.pdf}{Georgetown Journal of Law \& Public Policy}).

\section{Key Quotations}

\begin{quote}
``They had for more than a century before been regarded as beings of an inferior order, and altogether unfit to associate with the white race, either in social or political relations; and so far inferior, that they had no rights which the white man was bound to respect; and that the negro might justly and lawfully be reduced to slavery for his benefit.'' --- Chief Justice Roger B. Taney, Opinion of the Court (\href{https://constitutioncenter.org/the-constitution/supreme-court-case-library/dred-scott-v-sandford}{National Constitution Center})
\end{quote}

\begin{quote}
``And an act of Congress which deprives a citizen of the United States of his liberty or property merely because he came himself or brought his property into a particular Territory of the United States, and who had committed no offence against the laws, could hardly be dignified with the name of due process of law.'' --- Chief Justice Roger B. Taney, Opinion of the Court (\href{https://constitutioncenter.org/the-constitution/supreme-court-case-library/dred-scott-v-sandford}{National Constitution Center})
\end{quote}

\begin{quote}
``[T]he right of property in a slave is distinctly and expressly affirmed in the Constitution.'' --- Chief Justice Roger B. Taney, Opinion of the Court (\href{https://constitutioncenter.org/the-constitution/supreme-court-case-library/dred-scott-v-sandford}{National Constitution Center})
\end{quote}

\begin{quote}
``The truth is that citizenship, under the Constitution of the United States, is not dependent on the possession of any particular political or even of all civil rights, and any attempt so to define it must lead to error.'' --- Justice Benjamin R. Curtis, Dissenting Opinion (\href{https://famous-trials.com/dredscott/2545-scott-v-sandford-the-dissent}{Famous Trials})
\end{quote}

\begin{quote}
``And my opinion is, that, under the Constitution of the United States, every free person born on the soil of a State, who is a citizen of that State by force of its Constitution or laws, is also a citizen of the United States.'' --- Justice Benjamin R. Curtis, Dissenting Opinion (\href{https://constitutioncenter.org/the-constitution/supreme-court-case-library/dred-scott-v-sandford}{National Constitution Center})
\end{quote}

\end{document}
